\begin{afterword}
\summaryhead

The essay distinguishes wanting to investigate from wanting to have investigated. A person who questions their beliefs out of duty tends to consider a few objections, answer them, and stop, satisfied that they have done their duty. Their confidence ends up where it started. A curious person looks where their beliefs are most likely to change, and does not care which way they move.

Curiosity cannot be produced at will, so the author offers ways to work with duty. Watch for sparks of real interest and for possibilities you flinch from. Remember conservation of expected evidence: your expected future belief equals your present one, so each new piece of evidence should be able to move you either way. Ask whether you are rehearsing old arguments. Recite the Litany of Tarski (``If the box contains a diamond, I desire to believe that the box contains a diamond'') and imagine the advantages of the unwelcome answer. Guard any real uncertainty and let it grow into curiosity.

\responsehead

The essay has one distinction worth its length, between wanting to investigate and wanting to have investigated, and the author did not originate it; he borrows it through two quotations, one of them misattributed. Everything after the first paragraph is a set of remedies, and the remedies undercut the diagnosis.

The diagnosis says that duty-driven inquiry is hollow and that curiosity cannot be willed, any more than a cold foot can be willed warm. The remedies are then all acts of will: watch for sparks, meditate on a principle, recite a litany, imagine the benefits of the answer you fear. If the premise is true, the cure cannot work. If the cure works, the premise was false. The essay never notices the conflict.

Its one technical remedy is wrong. Conservation of expected evidence does not mean that belief should be ``evenly poised to shift in either direction,'' or that each new point has ``equal potential'' to move you up or down. It means only that the average of your possible future beliefs, weighted by their chances, equals your present belief. A confident person should expect small confirmations to be likely and large refutations to be rare. The author wrote exactly that two months earlier. Here he gives readers the wrong version and cites ``the laws of probability theory'' for it.

The post is also a closed circle of self-reference. It opens with a quotation from the author, quotes him again in a block introduced by the scriptural ``As it is written,'' borrows at least four more of his sentences without marks, links to his other posts for its key terms, and hands its positive account of curiosity to one of his fables. The Litany borrows a logician's name for what is a prayer. An essay about questioning your beliefs spends much of its length restating its author's.

It never shows anyone being curious. There is no worked case: no belief examined, no inquiry that went somewhere unexpected, no example of the ``flinching'' it warns against. The reader receives a mood and a vocabulary.

Its citations do not survive checking either. Twain credited the line to Caleb Winchester, and the post misquotes it. The Le Guin footnote cites a 2001 edition from an imprint first announced in 2013.

In short: an essay on curiosity that shows none, cures a failure of will with acts of will, and gets its one piece of mathematics wrong.
\end{afterword}
